\documentclass[11pt,a4paper]{article}

\usepackage[utf8]{inputenc}
\usepackage[T1]{fontenc}
\usepackage{lmodern}
\usepackage[margin=2.2cm]{geometry}
\usepackage{graphicx}
\usepackage{booktabs}
\usepackage{tabularx}
\usepackage{amsmath}
\usepackage{enumitem}
\usepackage{xurl}
\usepackage[hidelinks]{hyperref}

\setlength{\parindent}{0pt}
\setlength{\parskip}{6pt}
\setlist{itemsep=2pt, topsep=2pt}


\newcommand{\img}[1]{
  \begin{center}
    \includegraphics[width=\linewidth,height=0.8\textheight,keepaspectratio]{figures/#1}
  \end{center}}

\begin{document}

\begin{center}
  {\large Emanuele Girardello}\\[2pt]
  Matricola VR554844
\end{center}

%======================================================================
\section{Introduction}

\textbf{NOTE}\\
the original images serve as references
for evaluating the synthetic degradations, a clean reference may not be available in
real-world scenarios.

%======================================================================
\section{OBJECTIVE}
In this HomeWork we will implement a small pipeline for enhancing degraded images using classical image processing methods.
The objective is to analyze and evaluate the effects of the methods on the degraded images, understand and explain the reasons for such effects and justify the choices that are made.

the 4 methods we will implement are the following:

\begin{itemize}
  \item \textbf{baseline}: the degraded images, which are a result of a contrast reduction followed by a small amount of gaussian noise on the original images.
  \item \textbf{linear contrast stretching}: operating on both the BGR channels and on Luminance in the LAB format.
  \item \textbf{histogram\_equalization}: equalizing L on LAB.
  \item \textbf{bilateral filter + histogram equalization (LAB)}: the idea is to reduce the noise amplification of the histogram equalization by first applying a classical denoiser, this resulted in a reduction in noise, but also noticeable smoothing of the details compared to the single histogram equalization on L.
\end{itemize}

%======================================================================
\section{IMAGES AND SOURCES}

each image has been chosen based on a particular condition:

All images under unsplash license:\\
\url{https://unsplash.com/license}

low\_contrast: \url{https://unsplash.com/photos/a-blurry-photo-of-a-person-holding-a-cell-phone-RxQfhv2sBlU}\\
shadows: \url{https://unsplash.com/photos/a-wall-with-a-painting-on-it-cuiNuOWzilU}\\
poor\_details: \url{https://unsplash.com/photos/nighttime-cityscape-with-glowing-lanterns-and-lights-ZrDltpoqExE}

%======================================================================
\section{METHODS AND VISUAL COMPARISON}

Each figure shows, from left to right:
\begin{itemize}
  \item the original image (resized to 10\%)
  \item the degraded image (contrast reduction with $\alpha = 0.5$, $\beta = 64$, followed by Gaussian noise with $\sigma = 5$)
  \item the result of the corresponding method.
\end{itemize}

\subsection{Preprocessing and degradation parameters}

{\small
\begin{tabularx}{\linewidth}{@{}l X p{0.27\linewidth}@{}}
\toprule
Step & Code & Parameters \\
\midrule
Resize & \texttt{cv2.resize(img, None, fx=0.1, fy=0.1, interpolation=cv2.INTER\_AREA)} & factor 0.1, \texttt{INTER\_AREA} \\
Contrast reduction & \texttt{cv2.convertScaleAbs(img, alpha=0.5, beta=64)} & $\alpha = 0.5$, $\beta = 128\cdot(1 - \alpha) = 64$ \\
Gaussian noise & \texttt{cv2.randn(np.zeros(img.shape), (0), (noise\_sigma,) *3)} with \texttt{cv2.setRNGSeed(10)} & $\sigma = 5$ on all three channels, seed 10 \\
\bottomrule
\end{tabularx}
}

\texttt{convertScaleAbs} returns an 8-bit image; the noise is added in floating point, and the result is rounded and clipped to [0, 255] again.

\textbf{Resize to 10\%.} At 10\% scale folds and lights are still clearly visible, and execution is much faster. All window sizes below (bilateral filter, SSIM) are chosen with respect to the resized image.

\textbf{Interpolation: \texttt{INTER\_AREA} instead of the default \texttt{INTER\_LINEAR}.} \texttt The OpenCV documentation states that INTER\_AREA may be a preferred method for image decimation, as it gives moire'-free results. 

\textbf{Contrast reduction.}
\begin{itemize}
  \item $\alpha = 0.5$ halves every intensity difference. The contrast loss is strong, but every edge is still present, so the methods have something to recover.
  \item $\beta = 128\cdot(1 - \alpha) = 64$, Every pixel moves halfway toward 128 and the image is not darkened.
  This explains why the baseline PSNR depends so much on the image. low\_contrast is already close to mid-grey so the degradation changes it little.
  Shadows and poor\_details span almost 0--255 so the difference after the degradation is more noticeable.
  \item \texttt{cv2.convertScaleAbs} computes $|\alpha\cdot I + \beta|$ on B, G and R independently, then rounds and saturates to [0, 255].
\end{itemize}

\textbf{Gaussian noise.}
\begin{itemize}
  \item $\sigma = 5$, drawn independently for every pixel and channel.
  \item \textbf{Seed.} we set \texttt{cv2.setRNGSeed(10)}, so every run produces the same degraded images.
\end{itemize}

%======================================================================
\section{Linear contrast stretching}
we apply Linear contrast stretching, both Per-Channel for BGR and on Luminance for LAB
following the formula:
\[
\text{stretched} = \frac{\text{image} - f_{\min}}{f_{\max} - f_{\min}} \cdot (\text{new}_{\max} - \text{new}_{\min}) + \text{new}_{\min}
\]

\subsection{Parameters}

$f_{\min}$ and $f_{\max}$ are the lowest and highest values in the degraded image, computed separately for B, G and R (per-channel variant) or on L only (luminance variant). $\text{new}_{\min} = 0$ and $\text{new}_{\max} = 255$ use the whole 8-bit range.

Known limitation: this means $f_{\min}$ and $f_{\max}$ are most probably outliers generated by the added noise, a better solution would be to extract $f_{\min}$ and $f_{\max}$ as the 1 percentile and 99 percentile respectively.

%======================================================================
\section{Bilateral filter + histogram equalization (LAB)}

parameters for bilateral filter:
\begin{itemize}
  \item d=9, diameter of each pixel neighborhood.
  \item sigma\_color=30, standard deviation in the color space.
  \item sigma\_space=3, standard deviation in the coordinate space.
\end{itemize}

reasoning:

HE amplifies small differences, so a denoiser is applied first to reduce what HE will amplify. The bilateral filter replaces each pixel with a weighted mean of its neighbours. Each weight is the product of a spatial Gaussian (distance, sigma\_space) and a range Gaussian (color difference, sigma\_color), so flat regions are smoothed but strong edges are not.

\begin{itemize}
  \item \textbf{d = 9.}\\
  Determines the diameter of each pixel neighborhood that is used during filtering.
  \item \textbf{sigma\_color = 30.}\\
  A larger value of the parameter means that farther colors within the pixel neighborhood (see sigmaSpace) will be mixed together, resulting in larger areas of semi-equal color.
  The OpenCV documentation places useful values between 10 (little effect) and 150 (cartoon-like).
  \item \textbf{sigma\_space = 3.}\\
  A larger value of the parameter means that farther pixels will influence each other as long as their colors are close enough (see sigmaColor).
\end{itemize}

%======================================================================
\section{METRICS}
all the metrics are a result of the comparison of the method with the corresponding original image

\textbf{Metric parameters.}
\begin{itemize}
  \item \textbf{PSNR}: \texttt{cv2.PSNR(original, result)} with the default peak R = 255. The MSE is averaged over all three channels.
  \item \textbf{SSIM}: \texttt{skimage.metrics.structural\_similarity(original, result, channel\_axis=2)}\\
  library default parameters: $7\times7$ uniform window, K1 = 0.01, K2 = 0.03, data range 255 (inferred from uint8).\\
  SSIM is computed on B, G and R separately and then averaged.
\end{itemize}

\textbf{PSNR (dB)}

\begin{tabular}{@{}lrrr@{}}
\toprule
method & Shadows & low\_contrast & poor\_details \\
\midrule
degraded (baseline)       & 16.25 & 26.83 & 14.17 \\
stretch per channel (BGR) & 21.56 & 16.94 & 25.18 \\
stretch luminance (LAB)   & 22.46 & 21.21 & 26.51 \\
hist eq (LAB)             & 17.66 & 11.94 &  7.71 \\
bilateral + hist eq (LAB) & 17.76 & 12.01 &  6.38 \\
\bottomrule
\end{tabular}

\textbf{SSIM}

\begin{tabular}{@{}lrrr@{}}
\toprule
method & Shadows & low\_contrast & poor\_details \\
\midrule
degraded (baseline)       & 0.654 & 0.705 & 0.487 \\
stretch per channel (BGR) & 0.638 & 0.218 & 0.535 \\
stretch luminance (LAB)   & 0.710 & 0.297 & 0.669 \\
hist eq (LAB)             & 0.513 & 0.072 & 0.076 \\
bilateral + hist eq (LAB) & 0.728 & 0.401 & 0.070 \\
\bottomrule
\end{tabular}

%======================================================================
\section{GRID COMPARISON}

\textbf{low\_contrast}
\img{low_contrast_grid.jpg}

\textbf{poor\_details}
\img{poor_details_grid.jpg}

\textbf{shadows}
\img{Shadows_grid.jpg}

%======================================================================
\section{Numerical improvement vs.\ visual quality}

PSNR and SSIM are full reference metrics: they measure how close a result is to the original, not how good it looks.
Thus, a higher score does not always mean a better looking image.

\subsection{Shadows, bilateral filter + histogram equalization (LAB)}

This method has the highest SSIM of all on Shadows (0.728, against 0.710 for luminance stretching and 0.654 for the baseline), and its PSNR is above the baseline (17.76 against 16.25 dB).
But visually it is worse than luminance stretching:
\begin{itemize}
  \item the texture of the lower wall is flattened into low-detail patches.
  \item the light patches lose their warm color.
  \item the change in intensity between shadows and light is too drastic.
\end{itemize}

why SSIM rewards it:
SSIM compares the original and the modified image in small windows and, for each window, \textbf{multiplies} three terms:
\begin{itemize}
  \item \textbf{luminance}: local mean brightness similarity
  \item \textbf{contrast}: local amount of variation (standard deviation) similarity
  \item \textbf{structure}: correlation in pixel values.
\end{itemize}

In a flat window of the original (like a uniform piece of wall), any noise in the result adds variation that the original does not have, so the contrast and structure terms drop.\\
The bilateral filter removes that noise, so these windows score high again.

The bilateral filter gets high SSIM score only in the flat windows, where it removes noise, and gets a low SSIM score in the textured ones, where it removes detail.
the final SSIM score averages the two, so the gain on flat areas outweighs the loss of the texture we can actually see.

\subsection{Shadows, histogram equalization (LAB):}

PSNR improves over the baseline ($16.25 \rightarrow 17.66$ dB), but SSIM gets worse ($0.654 \rightarrow 0.513$). Visually the image presents noticeable issues:
\begin{itemize}
  \item the grain on the upper wall is heavily amplified.
  \item the light patches lose their warm color and turn almost white.
  \item once again, the change in intensity between shadows and light is too drastic.
\end{itemize}

why PSNR improves:\\
the degradation pulls every pixel halfway toward 128, which halves the spread in intensity, so the degraded image is wrong almost everywhere.
Equalization spreads the values out again, so most pixels move closer to their original value and the mean squared error drops.

why SSIM gets worse: equalization repairs the luminance and contrast terms, the two the degradation had damaged, but it amplifies the noise, so the structure term collapses. SSIM multiplies the three terms, so the collapse of the structure term drags the product down more than the gain in luminance and contrast lifts it.

\subsection{low\_contrast: the opposite case}

On low\_contrast, the baseline has the best PSNR and SSIM (26.83 dB / 0.705), yet it is the version where the folds are hardest to see.
Every other method lowers both metrics, because the original itself has low contrast: any enhancement moves the image away from the reference, even when it makes it easier to see the folds.

%======================================================================
\section{Qualitative analysis}

\subsection{Shadows: linear stretching on luminance (LAB)}

\begin{itemize}
  \item \textbf{Easier to see:} the edges of the shadows, the outline of the pole, the silhouette of the person and the texture of the wall. The dark wall keeps its neutral tone. This method has the best PSNR (22.46 dB), the best SSIM after bilateral + HE, and, after per-channel stretching, the colours closest to the original.
  \item \textbf{Trade-offs:} the noise is amplified and is visible as fine grain on the wall. The colours are slightly less saturated than in the original, because stretching only L does not restore the colors.
  \item \textbf{Other methods:}
  \begin{itemize}
    \item \emph{per-channel stretching} gives similar contrast, with warmer colors, but it also amplifies the noise.
    \item \emph{Histogram equalization} amplifies the grain on the wall and turns the light patches white. Applying the bilateral filter first removes the grain, but flattens the texture of the wall and does not fix the white lights.
  \end{itemize}
\end{itemize}

\subsection{low\_contrast: linear stretching on luminance (LAB)}

\begin{itemize}
  \item \textbf{Easier to see:} the folds and creases of the sheet, especially on the edges of the image, which are barely visible in the degraded image.
  \item \textbf{Trade-offs:} the noise is amplified as much as the folds so the image is grainy. The colours look more grey than in the original. Like before, stretching only L does not restore the colors.
  \item \textbf{Other methods:}
  \begin{itemize}
    \item \emph{per-channel stretching} produces unnatural colours, a strong blue and purple ``filter'' and a colored grain.
    \item \emph{Histogram equalization} shows the creases with the most contrast, but it amplifies the grain even more.
    \item \emph{bilateral filter + HE} cleans up the noise produced by HE, but details like small creases are lost, the borders also become dark.
  \end{itemize}
\end{itemize}

\subsection{poor\_details: linear stretching per channel (BGR)}

\begin{itemize}
  \item \textbf{Easier to see:} the lights in the image regain their contrast against the dark background, with colours closer to the original.
  \item \textbf{Trade-offs:} faint details on the buildings are lost. the image presents a colored grain that contrasts heavily with the dark background.
  \item \textbf{Other methods:}
  \begin{itemize}
    \item \emph{Luminance stretching (LAB)} gives almost the same image with less colored grain, it has the highest PSNR and SSIM (26.51 dB / 0.669), but the colours look faded: The degradation pulled the colors toward grey, and stretching only L restores the brightness but not the color.
    \item \emph{Histogram equalization (LAB)} fails on this image. Equalization tries to spread the pixels evenly over all Luminance levels. Here most of the image is dark, which in the degraded image results in a few L levels that differ only by the noise. Equalization spreads them over half of the output range, so the background turns into heavy salt-and-pepper noise. It highlights some details on the left, but at the cost of the rest of the image.
    \item \emph{Bilateral filter + histogram equalization (LAB)}: also fails. The filter reduces the noise, but the equalization mapping still amplifies the residual variations into a spotted black and white pattern.
  \end{itemize}
\end{itemize}

\vspace{1em}
link to the repository: \url{https://github.com/EmGira/Assignments_Computer_Vision/tree/main/Assignment1}

\end{document}
